\documentclass[10pt,a4paper]{amsart}
\usepackage[margin=19mm]{geometry}
\usepackage{amsmath,amssymb,mathtools}
\usepackage[scr=rsfs,cal=euler]{mathalfa}
\usepackage{xfrac}
\usepackage[unicode,hidelinks,bookmarks=false]{hyperref}
\newcommand{\ie}{i.e.\ }
\newcommand{\cf}{cf.\ }
\newcommand{\resp}{respectively\ }
\newcommand{\Subsection}{\subsection}
\providecommand{\nobreakdash}{\nobreak\hbox{-}}
\setlength{\emergencystretch}{2em}
\multlinegap0pt
\usepackage{bm}
\usepackage{bbm}
\usepackage{dsfont}
\usepackage{xspace}
\usepackage{cancel}
\usepackage{mathtools}
\usepackage{amsthm}
\makeatletter
\@ifundefined{theorem}{\newtheorem{theorem}{Theorem}}{}
\@ifundefined{claim}{\newtheorem{claim}[theorem]{Claim}}{}
\@ifundefined{fact}{\newtheorem{fact}[theorem]{Fact}}{}
\@ifundefined{lemma}{\newtheorem{lemma}[theorem]{Lemma}}{}
\makeatother
\def\BA{\mathsf{BA}}
\title[A note on adding isomorphisms and the pseudointersection num]{A note on adding isomorphisms and the pseudointersection number}
\author{Corey Bacal Switzer}
\date{Source: arXiv:2510.11155v1 (CC BY 4.0)}
\begin{document}
\maketitle

\noindent\emph{Source and licence.} A note on adding isomorphisms and the pseudointersection number. Version arXiv:2510.11155v1 (CC BY 4.0), distributed under the Creative Commons Attribution 4.0 International licence, as recorded in the arXiv licence metadata for this exact version and shown on the abstract page https://arxiv.org/abs/2510.11155v1. Full rights notice: licenses/arxiv-2510.11155-CC-BY-4.0.txt.\par\smallskip

\noindent\textit{Editorial scope.} Counted unit: the Lemma whose source label is \texttt{key} in the article (source0.tex lines 375--377, source label \texttt{key}), delivered together with the article's own complete original proof of it (source0.tex lines 379--402, one \texttt{proof} environment of 24 lines). No mathematics is written, repaired or re-derived: statement and proof are the article's own text line for line. Required context retained uncounted: none. Every definition, display and label of those lines is retained unchanged. Omitted from this package and not redistributed: 1--374, 378--378, 403--410. Nothing inside a retained excerpt is omitted or altered: every retained body is delivered exactly as the article's lines are written, so no omission receipt of any kind accompanies this package. Citations: the retained excerpts carry no citation command at all, so no bibliography entry of the article is transcribed and no reference is redistributed. Reference adaptation: none was needed and none was made; every reference inside the delivered unit resolves inside it, so no unresolved reference or citation is produced. Printed slips kept literal: no printed word, symbol, hypothesis or number of the counted statement or of its proof was altered. Layout and class: the article's own class is \texttt{amsart} with packages including bbm, xcolor, tikz, tikz-cd, setspace, babel, bussproofs, inputenc, csquotes, mathtools, mathrsfs, latexsym, enumitem, amsthm; the wrapper is the lane's pinned \texttt{amsart} editorial wrapper at 10pt on A4, which restates the theorem-like environments the retained text needs and adds the article's own macro definitions that the retained text needs (\texttt{\textbackslash BA}), with extra wrapper packages (bm, bbm, dsfont, xspace, cancel, mathtools). The delivered numbering is the unit's own. Assets: no retained span contains a figure environment, a graphics-inclusion command, a PDF-page inclusion, a table or a TikZ picture; no image file is copied into this package, the package declares no asset dependency and no image file is redistributed with it. Licence: version arXiv:2510.11155v1 of this article is distributed under the Creative Commons Attribution 4.0 International licence, as recorded in the arXiv licence metadata for this exact version and shown on the abstract page https://arxiv.org/abs/2510.11155v1. A full rights notice is delivered at licenses/arxiv-2510.11155-CC-BY-4.0.txt. Characters: every retained excerpt is pure ASCII, so the delivered engine, XeLaTeX with its default Latin Modern text font, reports no missing character.
\begin{lemma} \label{key}
    Assume $\BA$. Let $A$ and $B$ be $\aleph_1$-dense sets of reals. Let $F \subseteq \mathbb R \setminus A$ and $G \subseteq \mathbb R \setminus B$ be closed, nowhere dense sets and let $\varphi:F \to G$ be an order isomorphism between them. There is an isomorphism $\Phi:A \cup F \to B \cup G$ extending $\varphi$.
\end{lemma}

\begin{proof}
    Fix $F$, $G$, $A$, $B$ and $\varphi$ as in the statement of the lemma. We call an open interval $(a, b) \subseteq \mathbb R \setminus F$ {\em maximal} if it is not contained in any properly larger open interval disjoint from $F$. Let us denote by $M(F)$ the set of maximal open intervals disjoint from $F$ and similarly for $M(G)$ etc. The following can be checked easily.

    \begin{fact}
        The following all hold for $F$ (and also, for $G$).
        \begin{enumerate}
            \item Every $a \in \mathbb R \setminus F$ is contained in a single $I \in M(F)$.
            \item Every $I \in M(F)$ has the form $(a, b)$ for some $a < b \in F$.
            \item If $I, J \in M(F)$ are distinct then they are disjoint.
        \end{enumerate}
    \end{fact}

    The third item immediately implies that if $I , J \in M(F)$ then we can order them $I < J$ if and only if every $x \in I$ is less than every $y \in J$ and this is a total order on $M(F)$ (or $M(G)$ respectively). We will always treat $M(F)$ as a linear order with this ordering moving forward. The point is the following.

    \begin{claim}
        $\varphi:F \to G$ induces an order isomorphism $\hat{\varphi}:M(F) \cong M(G)$.
    \end{claim}

    \begin{proof}[Proof of Claim]
        If $I \in M(F)$ then by item (2) above $I = (a, b)$ for some $a < b \in F$ so we let $\hat{\varphi}(I) = (\varphi(a), \varphi(b))$. First let us see that this is well defined. If $I \in M(F)$ is as above the $\varphi(a) < \varphi(b) \in J$ and, since $\varphi$ is an order isomorphism and $F \cap (a, b)= \emptyset$ we must have that $\hat{\varphi}(I) \in M(G)$. Injectivity is immediate so let's check surjectivity. If $J \in M(F)$ then $J = (c, d)$ for some $c < d \in G$ and by an argument symmetric to the one proving well definedness we get that $(\varphi^{-1}(c), \varphi^{-1}(d)) \in M(F)$ and $J$ is the image of this interval under $\hat{\varphi}$. Thus $\hat{\varphi}$ is a well defined bijection between $M(F)$ and $M(G)$. Finally, if now $I < J \in M(F)$ then we must have that $a < b \leq c < d$ where $I = (a, b)$ and $J = (c, d)$. Consequently $\hat{\varphi}(I) < \hat{\varphi}(J)$ since $\varphi$ is assumed to be order preserving and hence we have $\varphi(a) < \varphi(b) \leq \varphi(c) < \varphi(d)$ so $\hat{\varphi}$ is order preserving.
    \end{proof}

Now, for each $I$ in $M(F)$ by $\BA$ there is an order isomorphism $\psi_I: I \cap A \cong \hat{\varphi}(I) \cap B$. We fix one for each such $I$. Finally let $\Phi:A \cup F \to B \cup G$ be the map defined simply as $\bigcup_{I \in M(F)}\psi_I \cup \varphi$. Let us check that this is the required isomorphism. It is nearly immediate from what has come before that this is a bijection, extends $\varphi$ and maps $A$ onto $B$. Finally since it is a union of order isomorphisms mapping ordered intervals to one another in order it is again order preserving. 
\end{proof}

\end{document}
